\section{Introduction}
\label{sec:introduction}

People increasingly oversee agent workflows that they neither program nor directly observe. A domain expert may decide what an agent should prioritize, approve its deployment, and bear the consequences of its actions, while an AI assistant or developer writes the instructions that guide those actions. We call the person who holds authority over the intended behavior and its release a \emph{workflow owner}. Their control is often mediated through reusable natural-language artifacts such as system prompts, custom instructions, templates, routines, and agent skills. We use \emph{skill} as an umbrella term for an artifact that organizes an agent's decisions and actions across a recurring family of tasks. Skills offer an editable control surface for studying how people understand and govern agent workflows.

A skill does not fully describe the workflow it will produce. The agent interprets it alongside the current task, surrounding instructions, available tools, intermediate results, and model variation. Each execution therefore enacts a context-specific sequence of decisions, tool calls, confirmation points, actions, and outcomes. We call this observable sequence the \emph{enacted workflow}. Even experienced maintainers who understand every sentence of a skill may struggle to anticipate this runtime behavior. Skill-mediated control consequently spans three objects: the behavior the owner intends, the editable artifact that expresses that intent, and the workflow the agent enacts.

Consider a researcher who relies on a travel-recovery skill. The skill checks affected commitments, retrieves feasible alternatives, compares arrival time and cost, applies airline and connection preferences, requests confirmation above a fare threshold, and rebooks. After a cancellation, the agent selects an \$80 same-airline flight arriving the following afternoon over a \$420 flight arriving that evening. This choice satisfies the stated cost and airline preferences, yet causes the researcher to miss a paper presentation the next morning. She immediately recognizes the outcome as unacceptable. The policy for similar cases remains unclear. The calendar-check instruction may have little influence on the later ranking rules; the cost rule may carry too much weight; the airline preference may apply too broadly; or the skill may need a priority rule for fixed commitments.

Several edits could repair this incident while encoding different policies for future travel. The skill could always prioritize the earliest arrival, prioritize arrival only when a fixed commitment is at risk, or relax cost and airline preferences for time-sensitive trips while retaining confirmation for expensive purchases. Each edit may select the desired flight in the motivating case. Their behavior diverges on routine trips, flexible meetings, costly alternatives, and journeys where airline or connection preferences should still matter. The failure identifies an objection without specifying the desired policy. This \emph{articulation gap} arises when an owner can reject an outcome before they can express the reusable policy that should replace it. A repair must therefore describe a relative behavioral change: what should change, what should remain stable, and which trade-offs remain unresolved.

Choosing a patch raises a second challenge. Natural-language rules may overlap, override one another, activate only in particular contexts, and yield variable executions. A local textual edit can therefore produce changes elsewhere in the workflow. A diff shows which words changed, while leaving the resulting decisions, tool calls, and outcomes uncertain. We call this the \emph{enactment gap} between an editable artifact and its runtime behavior. Owners see only a small sample of that behavior, which creates a related \emph{legibility gap} between observed executions and expectations for future use. Rerunning the motivating case may confirm a local fix, yet provides little evidence about transfer, regression, or stability.

Prior systems support important parts of this process. Prompt-authoring tools compare configurations and evaluate outputs against explicit criteria~\cite{arawjo2024chainforge,kim2024evallm}. Interactive debugging and what-if systems expose traces and alternative execution branches~\cite{epperson2025agdebugger,romerolauro2026rag}. Policy-authoring systems help practitioners articulate behavioral regions and expand test sets around ambiguous rules~\cite{lam2025policymaps,lee2026dataprompt}. Automated methods revise skills from execution evidence and rank candidates with task verifiers~\cite{liu2026skillrevise}. These systems increasingly support the translation of intent into prompts or policies and the comparison of artifacts or executions. Workflow owners still need to connect these activities when domain trade-offs resist a single verifier: articulate the intended change, understand how the skill shapes behavior across contexts, and judge whether a concrete revision is ready for release.

We ask: \emph{How can interactive systems help workflow owners understand and govern agent workflows that they control indirectly through reusable natural-language skills?} We study repair episodes, where the relationship among intent, artifact, and enacted behavior becomes consequential. We frame this activity as \emph{behavioral patch review}. Like software patch review, it asks a person with decision authority to inspect a proposed change before release. LLM agents make the object of review less direct: repeated executions must reveal the behavioral diff, while judgments over concrete cases help the owner articulate the intended diff. We define a well-scoped repair as the smallest behavioral change justified by the owner's recorded intent and the available evidence.

We instantiate this interaction model in \systemname{}. First, the owner builds a versioned behavioral scope from concrete cases, marking behavior to \emph{change}, \emph{preserve}, or leave \emph{unresolved}. Second, controlled interventions identify instructions in the original skill that influence these commitments. \systemname{} uses this evidence to enrich a familiar line-based diff, linking each candidate edit to its intended purpose and relevant executions. Third, the system runs the complete original and revised skills under matched conditions and compares the owner's \emph{intended behavioral delta} with the observed \emph{enacted behavioral delta}. Boundary judgments recorded before candidate outcomes provide additional checks on the reach of the repair. The owner can then revise the scope or skill, gather more evidence, defer the repair, or release the new version.

We investigate behavioral patch review at three levels. A formative study of real repair episodes examines how workflow owners relate intended behavior, reusable instructions, and enacted workflows. A technical evaluation assesses evaluator fidelity and the stability and discriminative validity of controlled interventions. A controlled study measures users' predictions of held-out skill behavior, identification of repair-relevant instructions, alignment between intended and enacted repair scope, regression detection, and release decisions.
% TODO before submission: Insert a concise, verified findings paragraph here. Lead with the
% primary human outcome, then scope mismatch, release/abstention quality, and the most
% important technical boundary. Do not reuse the synthetic values from the planning draft.

This work makes three contributions:
\begin{enumerate}
    \item \textbf{An empirical account of skill repair by workflow owners.} Through real repair episodes, we characterize how owners connect intended behavior, reusable skill instructions, and context-specific executions. The findings reveal recurring gaps in articulating a repair, anticipating how instructions shape behavior, and judging whether a revision is ready for reuse.
    \item \textbf{Behavioral patch review, an interaction model for governing changes to reusable agent skills.} The workflow connects a versioned behavioral scope to an evidence-augmented skill diff and matched executions of the complete revision.
    \item \textbf{Technical and human evidence for behavioral patch review.} We evaluate the evidence underlying the interface and study how the approach affects skill understanding, repair-scope alignment, regression detection, and release decisions.
\end{enumerate}

Together, these contributions position skills as editable intermediaries between user intent and agent behavior. They show how interfaces for maintaining reusable agents can make behavioral change reviewable by the people who own its consequences, including those who do not directly engineer the workflow.
